%==============================================================
% ΚΕΦΑΛΑΙΟ 2 — ΘΕΩΡΗΤΙΚΟ & ΚΑΝΟΝΙΣΤΙΚΟ ΠΛΑΙΣΙΟ
%==============================================================
\chapter{Θεωρητικό και Κανονιστικό Πλαίσιο}
\label{ch:theory}

Το παρόν κεφάλαιο παρουσιάζει το θεωρητικό και κανονιστικό υπόβαθρο
που διέπει την επικύρωση μοντέλων μηχανικής μάθησης στον πιστωτικό
κίνδυνο. Αρχικά εξετάζεται η έννοια του πιστωτικού κινδύνου και η
ιστορική εξέλιξη των μεθόδων μέτρησής του, ακολουθεί η παρουσίαση
του κανονιστικού πλαισίου επικύρωσης, και στη συνέχεια αναλύονται
οι τεχνικές ερμηνευσιμότητας, ανίχνευσης μεροληψίας και
παρακολούθησης μοντέλων.

%==============================================================
\section{Πιστωτικός Κίνδυνος και Credit Scoring}
\label{sec:credit_risk}
%==============================================================

%--------------------------------------------------------------
\subsection{Ορισμός και Ποσοτικοποίηση του Πιστωτικού Κινδύνου}
\label{subsec:credit_risk_def}
%--------------------------------------------------------------

Κάθε φορά που ένα πιστωτικό ίδρυμα χορηγεί δάνειο, αναλαμβάνει τον
κίνδυνο ο δανειολήπτης να αδυνατεί να ανταποκριθεί στις υποχρεώσεις
του. Ο κίνδυνος αυτός ορίζεται ως \textbf{πιστωτικός κίνδυνος} και
αποτελεί έναν από τους σημαντικότερους κινδύνους που αντιμετωπίζει ο
τραπεζικός τομέας \citep{basel2006}.

Για την ποσοτικοποίησή του, το εποπτικό πλαίσιο της Βασιλείας
(Basel~II/III) εισήγαγε τρεις θεμελιώδεις παραμέτρους:

\begin{itemize}
  \item \textbf{Πιθανότητα Αθέτησης} (\textit{Probability of Default,
  PD}): η πιθανότητα ο δανειολήπτης να κηρυχθεί σε αδυναμία πληρωμής
  εντός ενός δεδομένου χρονικού ορίζοντα, συνήθως ενός έτους.

  \item \textbf{Έκθεση κατά την Αθέτηση} (\textit{Exposure at Default,
  EAD}): το ανεξόφλητο υπόλοιπο της οφειλής τη στιγμή της αθέτησης
  πληρωμής.

  \item \textbf{Ζημία σε Περίπτωση Αθέτησης} (\textit{Loss Given
  Default, LGD}): το ποσοστό του EAD που δεν ανακτάται μετά από
  ενέργειες όπως ο πλειστηριασμός εξασφαλίσεων ή η νομική
  διεκδίκηση.
\end{itemize}

Ο συνδυασμός των τριών αυτών παραμέτρων δίνει την
\textbf{Αναμενόμενη Ζημία} (\textit{Expected Loss, EL}):

\begin{equation}
  \text{EL} = \text{PD} \times \text{LGD} \times \text{EAD}
  \label{eq:expected_loss}
\end{equation}

Η δυνατότητα ποσοτικοποίησης του πιστωτικού κινδύνου μέσω της
εξίσωσης~\eqref{eq:expected_loss} είναι κρίσιμη για τα πιστωτικά
ιδρύματα. Τους επιτρέπει να λαμβάνουν τεκμηριωμένες αποφάσεις
χορήγησης δανείων, να διαμορφώνουν κατάλληλο επιτόκιο ανάλογα με το
προφίλ κινδύνου κάθε δανειολήπτη, και να διατηρούν επαρκή κεφαλαιακά
αποθέματα σύμφωνα με τις εποπτικές απαιτήσεις \citep{mcneil2015}.

%--------------------------------------------------------------
\subsection{Ιστορική Εξέλιξη των Μεθόδων Credit Scoring}
\label{subsec:credit_scoring_history}
%--------------------------------------------------------------

Οι τεχνικές για τον υπολογισμό του πιστωτικού κινδύνου μπορούν να 
κατηγοριοποιηθούν σε τρεις βασικές εποχές, όπου η κάθε μια έχει 
καλύτερη ακρίβεια από την προηγούμενη, αλλά με διαφορετικά 
μειονεκτήματα και ανάγκες για ομαλή εφαρμογή. Η πρώτη εποχή, μέχρι 
το 1980, χαρακτηρίζεται από την κρίση των ειδικών 
(\textit{Expert Judgment}), όπου ένας έμπειρος αναλυτής διάβαζε τα 
δεδομένα και έκανε μια πρόβλεψη βασισμένη στην εμπειρία του. Στη 
συνέχεια, κατά την περίοδο 1980--2010, κυριάρχησε η χρήση τεχνικών 
όπως τα \textit{Statistical Scorecards} και η Logistic Regression, με 
κύριο χαρακτηριστικό την επαναληψιμότητα και τη μεγαλύτερη ακρίβεια 
σε σχέση με την υποκειμενική κρίση των ειδικών \citep{thomas2000}. 
Τέλος, από τη δεκαετία του 2010 και έπειτα, η μηχανική μάθηση έχει 
αποκτήσει κυρίαρχο ρόλο. Σε αντίθεση με τις τεχνικές της 
προηγούμενης εποχής, τα μοντέλα μηχανικής μάθησης μπορούν να κάνουν 
προβλέψεις χωρίς να υποθέτουν γραμμικότητα στη σχέση των 
μεταβλητών, επιτυγχάνοντας έτσι υψηλότερη προβλεπτική ακρίβεια, 
αλλά μειώνοντας παράλληλα την ερμηνευσιμότητα \citep{hastie2009}. 
Η ανάγκη αντιμετώπισης αυτής της έλλειψης ερμηνευσιμότητας, καθώς 
και η παρακολούθηση της συμπεριφοράς των μοντέλων, αποτελούν το 
κεντρικό ερώτημα της παρούσας εργασίας.

%--------------------------------------------------------------
\subsection{Πιστωτική Βαθμολόγηση έναντι Μοντελοποίησης Κινδύνου}
\label{subsec:scoring_vs_modeling}
%--------------------------------------------------------------
Η μελέτη του πιστωτικού κινδύνου διακρίνεται στην πιστωτική 
βαθμολόγηση (\textit{credit scoring}) και τη μοντελοποίηση 
πιστωτικού κινδύνου (\textit{credit risk modeling}) \citep{thomas2000}. 
Οι δύο αυτές κατηγορίες έχουν σημαντικές διαφορές, με την πρώτη να 
αφορά αποκλειστικά την πιθανότητα αθέτησης (\textit{Probability of 
Default, PD}) σε ένα μεμονωμένο δάνειο, ενώ η δεύτερη αφορά 
ολόκληρο το χαρτοφυλάκιο μιας τράπεζας, το οποίο αποτελείται από 
πολλαπλά δάνεια και απαιτεί εκτίμηση συνολικών ζημιών, stress testing 
και κεφαλαιακές απαιτήσεις \citep{mcneil2015}. Η παρούσα εργασία 
εστιάζει στην πιστωτική βαθμολόγηση και συγκεκριμένα στην αξιολόγηση 
μοντέλων μηχανικής μάθησης για την πρόβλεψη της PD σε ατομικό επίπεδο.

%==============================================================
\section{Μοντέλα Μηχανικής Μάθησης στον Πιστωτικό Κίνδυνο}
\label{sec:ml_models}
%==============================================================

Τα ML μοντέλα παρέχουν υψηλότερη ακρίβεια σε σχέση με τα παραδοσιακά 
διαθέσιμα εργαλεία και αυτός είναι ο κύριος λόγος για τον οποίο η 
χρήση τους γίνεται ολοένα και πιο διαδεδομένη \citep{hastie2009}. 
Οι δύο βασικές κατηγορίες αυτών των μοντέλων είναι το Gradient 
Boosting και τα νευρωνικά δίκτυα. Η πρώτη κατηγορία αξιοποιεί τη 
θεωρία των δέντρων αποφάσεων (\textit{decision trees}) και προσφέρει 
μερική ερμηνευσιμότητα, ενώ η δεύτερη στηρίζεται στα τεχνητά 
νευρωνικά δίκτυα με ελάχιστη ερμηνευσιμότητα. Παρόλο που η Logistic 
Regression υστερεί σε ακρίβεια έναντι των παραπάνω, αποτελεί ορόσημο 
στον τομέα και χρησιμοποιείται ευρέως ως σημείο αναφοράς 
\citep{siddiqi2006}. Στην παρούσα εργασία γίνεται χρήση και των τριών 
προσεγγίσεων --- της Logistic Regression ως παραδοσιακό μοντέλο 
αναφοράς, του XGBoost ως αντιπροσωπευτικό μοντέλο Gradient Boosting, 
και ενός Feed-Forward Neural Network ως παράδειγμα αυξημένης 
πολυπλοκότητας.

%--------------------------------------------------------------
\subsection{Logistic Regression}
\label{subsec:logistic_regression}
%--------------------------------------------------------------

Η Logistic Regression είναι τυπικά ένας παραδοσιακός αλγόριθμος 
μηχανικής μάθησης ο οποίος βασίζεται στη Γραμμική Παλινδρόμηση 
(\textit{Linear Regression}), της οποίας το αποτέλεσμα 
μετασχηματίζεται στο διάστημα $(0,1)$ μέσω της συνάρτησης sigmoid:
\begin{equation}
  \sigma(z) = \frac{1}{1 + e^{-z}}
  \label{eq:sigmoid}
\end{equation}
Όταν το $z$ είναι μεγάλο η συνάρτηση πλησιάζει στο 1, που 
μεταφράζεται σε υψηλή πιθανότητα αθέτησης, ενώ όταν είναι μικρό 
πλησιάζει στο 0, δηλαδή χαμηλή πιθανότητα αθέτησης. Για την εύρεση 
του $z$ ακολουθείται ο γραμμικός συνδυασμός:
\begin{equation}
  z = \beta_0 + \beta_1 x_1 + \beta_2 x_2 + \ldots + \beta_n x_n
  \label{eq:logistic_z}
\end{equation}
όπου οι συντελεστές $\beta$ εκτιμώνται από τα δεδομένα 
\citep{hastie2009}. Με αυτόν τον τρόπο καθίσταται εφικτή η ερμηνεία 
του αποτελέσματος, καθώς φαίνεται άμεσα ποιοί συντελεστές επηρέασαν 
περισσότερο την τελική πρόβλεψη. Σημαντικό όριο της τεχνικής είναι 
ότι κάθε χαρακτηριστικό θεωρείται γραμμικό ως προς την τελική πρόβλεψη, 
με αποτέλεσμα να μην μπορεί να αναπαραστήσει σύνθετες 
αλληλεπιδράσεις μεταξύ χαρακτηριστικών, όπως ``υψηλό εισόδημα 
\textsc{και} νέα ηλικία \textsc{και} καμία πιστωτική ιστορία 
$\rightarrow$ υψηλός κίνδυνος'' \citep{thomas2000}. Για τους λόγους 
αυτούς χρησιμοποιείται στην παρούσα εργασία ως μοντέλο αναφοράς 
(\textit{baseline}).

%--------------------------------------------------------------
\subsection{Gradient Boosting}
\label{subsec:gradient_boosting}
%--------------------------------------------------------------

Το Gradient Boosting είναι ένα σύγχρονο εργαλείο μηχανικής μάθησης 
το οποίο βασίζεται στη θεωρία των δέντρων αποφάσεων 
(\textit{Decision Trees}) \citep{hastie2009}. Τα δέντρα αποφάσεων 
είναι δομές που με απλό τρόπο, μέσω συγκρίσεων ανάλογα με το βάθος 
τους, οδηγούν σε ένα αποτέλεσμα. Η ιδέα του Boosting χρησιμοποιεί 
πολλαπλά δέντρα συνήθως μικρού βάθους, τα οποία μαθαίνουν από τα 
λάθη των προηγούμενων και συνδυάζουν τα αποτελέσματά τους. Αυτό 
αποτελεί τη βάση του αλγορίθμου Gradient Boosting, ο οποίος έχει 
πολλαπλές υλοποιήσεις. Στην παρούσα εργασία χρησιμοποιείται η 
υλοποίηση XGBoost \citep{chen2016xgboost}, η οποία επιλέχθηκε λόγω 
της υψηλής ταχύτητάς της, του ενσωματωμένου μηχανισμού regularization 
που αποτρέπει την υπερπροσαρμογή, και της ευρείας χρήσης της στον 
τομέα του πιστωτικού κινδύνου. Ο μεγάλος αριθμός των δέντρων και η 
αλληλοσυσχέτισή τους καθιστά αδύνατη τη σαφή ερμηνεία των 
αποτελεσμάτων. Για τον λόγο αυτόν κρίνεται αναγκαία η χρήση της 
μεθόδου SHAP \citep{lundberg2017shap}, μέσω της οποίας καθίσταται 
δυνατή η ποσοτική ερμηνεία της συνεισφοράς κάθε χαρακτηριστικού 
στο αποτέλεσμα.

%--------------------------------------------------------------
\subsection{Feed-Forward Neural Network}
\label{subsec:neural_network}

Το Feed-Forward Neural Network (Νευρωνικό Δίκτυο Πρόσθιας Τροφοδότησης, FFNN)
αποτελείται από στρώματα νευρώνων. Η λειτουργία κάθε νευρόνα γίνεται εύκολα
κατανοητή μέσω του βάρους του, αλλά το σύνολό τους δημιουργεί ένα περίπλοκο
σύστημα αδύνατο να κατανοηθεί αναλυτικά. Για αυτόν τον λόγο η τεχνική αυτή
χαρακτηρίζεται ως μαύρο κουτί (black box). Με την χρήση τεχνικών post-hoc XAI
μπορούμε να προσεγγίσουμε την λογική πίσω από τους νευρώνες, χωρίς όμως
απόλυτα σαφή αποτελέσματα. Η έλλειψη απόλυτης ερμηνευσιμότητας και η
πολυπλοκότητα οδηγούν σε πιο απαιτητικό πλαίσιο επικύρωσης (validation
framework). Στην εργασία γίνεται χρήση του FFNN ως σημείου αναφοράς
πολυπλοκότητας ώστε να αναδειχθεί πως η μειωμένη ερμηνευσιμότητα οδηγεί σε
μεγαλύτερες απαιτήσεις επικύρωσης χωρίς αντίστοιχο όφελος σε απόδοση
έναντι του XGBoost.
%==============================================================
\section{Επικύρωση Μοντέλων (Model Validation)}
\label{sec:model_validation}
%==============================================================

Η επικύρωση μοντέλου (Model Validation) είναι η διαδικασία ελέγχου
του μοντέλου στην ικανότητά του να παράγει ικανοποιητικά αποτελέσματα
μετά την εκπαίδευσή του. Το σύστημα μαθαίνει να παίρνει αποφάσεις με
παραδείγματα από μια συγκεκριμένη βάση δεδομένων. Καμία βάση δεδομένων
δεν γίνεται να περιέχει όλα τα δυνατά σενάρια που θα κληθεί να
αντιμετωπίσει το μοντέλο, άρα θα πρέπει να αποδείξει την ικανότητά
του να γενικεύει σωστά σε νέα δεδομένα. Συστήματα μηχανικής μάθησης
που δεν έχουν αυτή την ικανότητα μπορούν να κοστίσουν στο τραπεζικό
σύστημα αλλά και στο κοινωνικό σύνολο με τις λανθασμένες εκτιμήσεις
τους.

Η επικύρωση των μοντέλων είναι νομική υποχρέωση. Οι τράπεζες που
χρησιμοποιούν εσωτερικά μοντέλα υποχρεούνται μέσω του πλαισίου της
Βασιλείας II/III (Basel II/III) να αποδεικνύουν την εγκυρότητα και
αξιοπιστία τους \citep{basel2006, basel2017}. Με μικρές διαφορές,
τόσο το SR Letter 11-7 της Federal Reserve των ΗΠΑ όσο και οι EBA
Guidelines στην Ευρώπη επιβάλλουν την επικύρωση, παρακολούθηση αλλά
και εξήγηση και δικαιολόγηση των τεχνικών μηχανικής μάθησης
\citep{sr117, eba2023}.

Το κύριο παραδοσιακό μοντέλο, η Logistic Regression, έχει το βασικό
πλεονέκτημα ότι είναι μια απόλυτα ερμηνεύσιμη διαδικασία. Εάν
απορριφθεί κάποιο δάνειο, κοιτώντας τους ανάλογους συντελεστές
μπορεί να καταλάβει κανείς γιατί πάρθηκε αυτή η απόφαση. Η διαδικασία
αυτή δυσκολεύει με τα σύγχρονα μοντέλα ML. Αλγόριθμοι όπως το XGBoost
και το FFNN χρειάζονται σύνθετες αναλύσεις SHAP για να κατανοήσουμε
τις μεταβλητές που οδηγούν στο εκάστοτε αποτέλεσμα \citep{lundberg2017shap}.
Απαιτείται επίσης έλεγχος για μεροληψία (bias) εναντίον κοινωνικών
ομάδων αλλά και συνεχής παρακολούθηση (monitoring) για επικύρωση της
ομαλής τους λειτουργίας καθώς τα δεδομένα μεταβάλλονται. Οι διαφορές
αυτές στις απαιτήσεις ανάμεσα στα παραδοσιακά και σύγχρονα μοντέλα
αποτελούν το κεντρικό ερώτημα της εργασίας.

%==============================================================
\section{Ερμηνευσιμότητα (Interpretability / XAI)}
\label{sec:interpretability}
%==============================================================
Ερμηνευσιμότητα (Interpretability) είναι η ιδιότητα του μοντέλου να
μπορεί να γίνει κατανοητό με ανθρώπινους όρους. Η ιδιότητα αυτή είναι
νομικά απαραίτητη και χρήσιμη τόσο στους πελάτες όσο και στους
αναλυτές που παρακολουθούν το μοντέλο. Υπάρχουν δύο κατηγορίες
ερμηνευσιμότητας, η εγγενής (intrinsic) και η post-hoc. Η πρώτη
περιγράφει μοντέλα που από τη φύση τους είναι ερμηνεύσιμα, όπως είναι
η Logistic Regression, ενώ η δεύτερη περιγράφει μοντέλα που απαιτούν
τεχνικές όπως το SHAP και το LIME για την ερμηνεία τους
\citep{molnar2022}.

Πιο συγκεκριμένα, το SHAP (SHapley Additive exPlanations) προέρχεται
από τον κλάδο της θεωρίας παιγνίων των μαθηματικών και μετράει τη συνεισφορά κάθε
παίκτη σε ένα συνεργατικό παιχνίδι \citep{lundberg2017shap}. Στην
περίπτωση των αλγορίθμων μηχανικής μάθησης οι παίκτες είναι τα
χαρακτηριστικά του εν δυνάμει δανειολήπτη, ενώ η τελική συνεισφορά
είναι η πρόβλεψη. Ως τελικό αποτέλεσμα παίρνουμε συντελεστές που
φανερώνουν τη βαρύτητα κάθε χαρακτηριστικού. Καταλαβαίνουμε έτσι
ποια χαρακτηριστικά θεωρεί σημαντικά το μοντέλο μας και ποια
χαρακτηριστικά του υποψήφιου δανειολήπτη επηρέασαν αρνητικά την
αίτησή του.

Σε αντίθεση με το SHAP, το LIME (Local Interpretable Model-agnostic
Explanations) εστιάζει μεμονωμένα σε κάθε πρόβλεψη
\citep{ribeiro2016lime}. Λόγω αυτού είναι ταχύτερο από το SHAP αλλά
περισσότερο ασταθές. Για τον λόγο αυτό χαίρει μικρότερης αποδοχής
στον κλάδο και έτσι στην παρούσα εργασία θα εστιάσουμε στην πρώτη
μέθοδο.

%==============================================================
\section{Μεροληψία και Δικαιοσύνη (Bias \& Fairness)}
\label{sec:bias}
%==============================================================

Μεροληψία (bias) στο πλαίσιο της μηχανικής μάθησης είναι η τάση του
μοντέλου να παίρνει λάθος αποφάσεις για συγκεκριμένες ομάδες του
πληθυσμού. Γενικά οι ομάδες αυτές μπορούν να αφορούν το φύλο, την
εθνικότητα ή την ηλικία. Στην παρούσα εργασία όμως, λόγω της βάσης
δεδομένων που χρησιμοποιείται, η οποία δεν περιέχει τέτοια δεδομένα,
οι ομάδες αυτές ορίζονται με βάση χαρακτηριστικά όπως το μέγεθος του
πιστωτικού ιστορικού. Εάν το μοντέλο κάνει διακρίσεις υπέρ των
πελατών με μεγάλο πιστωτικό ιστορικό, τότε η αξιοπιστία του θα φθίνει
και θα οδηγήσει σε λάθος αποφάσεις. Τελικά δεν στοχεύουμε στο να μην
λαμβάνει καθόλου υπόψη αυτή τη μεταβλητή το μοντέλο μας, αλλά να της
δίνει τη βαρύτητα που της αναλογεί χωρίς υπερβολές.

Για την μέτρηση της μεροληψίας μοντέλου θα γίνει χρήση των μετρικών
της Δημογραφικής Ισότητας (Demographic Parity) και της Ανισομερούς
Επίπτωσης (Disparate Impact) \citep{barocas2019fairness}. Η
Δημογραφική Ισότητα ελέγχει εάν δίνονται δάνεια με την ίδια συχνότητα
σε όλες τις ομάδες, ενώ η Ανισομερής Επίπτωση ελέγχει εάν
λαμβάνονται δυσανάλογα αρνητικές αποφάσεις απέναντι σε ένα μέρος του
συνόλου. Πιο αναλυτικά, η πρώτη θέτει ως ιδανικό το ίδιο ποσοστό
έγκρισης για όλες τις ομάδες. Η δεύτερη υπολογίζει τον λόγο έγκρισης
μεταξύ των ομάδων και θέτει το αποδεκτό όριο στο 0,8. Οι μετρικές
αυτές θα εφαρμοστούν και στα τρία μοντέλα: Logistic Regression,
XGBoost και FFNN.

Η προσπάθεια ανίχνευσης μεροληψίας θα γίνει στα μοντέλα Logistic
Regression, XGBoost και FFNN, με το τρίτο να έχει επιδεικτικό ρόλο
στη σχέση μεταξύ πολυπλοκότητας και ανίχνευσης αλλά και ερμηνείας
της μεροληψίας, ενώ θα χρησιμοποιηθούν υλοποιήσεις από τη βιβλιοθήκη
Fairlearn \citep{bird2020fairlearn}. Η διάγνωση bias από μόνη της δεν
είναι τόσο χρήσιμη και πρέπει να συνοδεύεται από ερμηνεία. Η αιτία
δημιουργίας μεροληψίας μπορεί να βρίσκεται τόσο στο ίδιο το μοντέλο
και τις διαδικασίες εκμάθησής του όσο και στην ίδια τη βάση δεδομένων.
Το γεγονός αυτό καθιστά σαφές ότι όσο πιο σύνθετο είναι ένα μοντέλο,
τόσο πιο απαιτητική είναι όλη η διαδικασία.
%==============================================================
\section{Παρακολούθηση Μοντέλων (Model Monitoring)}
\label{sec:monitoring}
%==============================================================

Η εκπαίδευση του εκάστοτε συστήματος γίνεται ώστε αυτό να μπορεί να
λειτουργεί στις τρέχουσες οικονομικές συνθήκες και να είναι αποδοτικό
σε αυτές. Με το πέρασμα του χρόνου παρατηρούνται μεταβολές σε αυτές
τις συνθήκες, οπότε η παραγωγή προβλέψεων χωρίς παρακολούθηση γεννά
κινδύνους. Το Model Monitoring είναι η διαδικασία παρακολούθησης του
μοντέλου μετά το πέρας της ανάπτυξής του.

Τα μοντέλα τεχνητής νοημοσύνης μετά την εκπαίδευσή τους συνεχίζουν
να λειτουργούν με τον ίδιο τρόπο, ο οποίος όμως δεν παράγει σε όλες
τις συνθήκες ορθά αποτελέσματα. Τα δεδομένα εισόδου και οι
συσχετίσεις μεταξύ αυτών και του αποτελέσματος είναι πιθανόν να
μετατοπιστούν, δημιουργώντας ανάγκη για αλλαγή στη διαχείριση των
δεδομένων από το μοντέλο. Η αλλαγή της κατανομής των δεδομένων
εισόδου ονομάζεται Data Drift και στην περίπτωσή μας μπορεί να συμβεί
ως αποτέλεσμα μιας οικονομικής ύφεσης \citep{widmer1996learning}. Η
δεύτερη περίπτωση, η αλλαγή των συσχετίσεων, ονομάζεται Concept Drift
και έχει ως αποτέλεσμα οι ίδιες μεταβλητές που παρήγαγαν ορθό
αποτέλεσμα πλέον να προκαλούν σφάλμα. Μπορεί να εξηγηθεί απλοϊκά ως
αλλαγή των κανόνων του παιχνιδιού.

Η παρακολούθηση γίνεται μέσω των μετρικών PSI (Population Stability
Index) και CSI (Characteristic Stability Index) \citep{eba2023}. Η PSI
εστιάζει σε αλλαγές της κατανομής των δεδομένων και η CSI σε κάθε
μεταβλητή ξεχωριστά. Τα όριά τους είναι συγκεκριμένα και φανερώνουν
εάν το μοντέλο είναι σταθερό ή εάν χρειάζεται επανεκπαίδευση.

Η εργασία εφαρμόζει τις παραπάνω τεχνικές monitoring αναλυτικά στον
αλγόριθμο XGBoost, ενώ για το FFNN η παρακολούθηση εστιάζει
περισσότερο στην επίδειξη της αυξημένης δυσκολίας λόγω της αισθητά
μεγαλύτερης πολυπλοκότητάς του, όπου τα περισσότερα στρώματα
νευρώνων αφήνουν λιγότερο χώρο για διαφάνεια και ακριβή
παρακολούθηση.